\scanpage{12}
126\hfill J.-P. \textsc{Serre}

 où $E'$ (resp. $E''$) est un $S'$-module (resp. $S''$-module) absolument simple. Soit $\widetilde\Psi$ la composante neutre de l'image réciproque de $\Psi$ dans $S'\times S''$. C'est un tore de dimension 1, donc isomorphe au groupe $\mathbf G_m$. Choisissons un isomorphisme $\sigma:\widetilde\Psi\to\mathbf G_m$. Notons respectivement $f$, $f_0$, $f_1$, $f'$, $f''$ les caractères des représentations $E$, $E_0$, $E_1$, $E'$, $E''$ de $\widetilde\Psi$. Ce sont des polynômes à coefficients entiers positifs en $\sigma$ et $\sigma^{-1}$. On a
\[
f'\cdot f''=f=f_0+f_1
\]
et
\[
f_0=n_0\sigma^a,\quad f_1=n_1\sigma^b,\quad\text{avec}\quad a,b\in\mathbf Z.
\]

On a $a\ne b$; sinon, en effet, $\Psi$ opèrerait sur $E$ par homothéties. De plus:
\[
f=n_0\sigma^a+n_1\sigma^b.
\]

Mais on vérifie tout de suite qu'un tel binôme ne peut se décomposer en produit de deux polynômes à coefficients entiers $\geq0$ que de façon triviale, l'un des facteurs étant un monôme $c\sigma^d$. L'entier $c$ doit diviser $n_0$ et $n_1$, donc doit être égal à 1 vu l'hypothèse $\operatorname{pgcd}(n_0,n_1)=1$. On a donc, par exemple, $f'=\sigma^d$, d'où $\dim.E'=1$. Mais c'est absurde, car l'image de $S'$ dans $S$ opèrerait trivialement sur $E$, contrairement au fait que $E$ est un $S$-module fidèle. Le groupe $S$ est donc bien un groupe \textit{simple}.

4) Soit $h=n_0+n_1$ la dimension de $E$, et soit $\mu_h$ le groupe des racines $h$-èmes de l'unité. On a $\lambda^{n_1}=\lambda^{-n_0}$ si $\lambda\in\mu_h$. L'homomorphisme $\psi$ de 2) transforme donc $\lambda\in\mu_h$ en une \textit{homothétie}. Comme $n_0$ et $n_1$ sont premiers entre eux, $\psi$ est injectif. On en conclut finalement que $\Psi$, donc \textit{a fortiori} $S$, \textit{contient le groupe $\mu_h$}, identifié à un sous-groupe du groupe des homothéties. En particulier, \textit{$\mu_h$ est contenu dans le centre de $S$}.

5) Vu 3) et 4), il ne nous reste plus qu'à démontrer le lemme suivant:

\textbf{Lemme 2.} -- \textit{Soit $S$ un groupe algébrique simple, et soit $\varrho:S\to\mathbf{GL}_E$ une représentation linéaire non triviale de $S$, de dimension $h$. Si le centre de $S$ est d'ordre multiple de $h$, $\varrho$ est un isomorphisme de $S$ sur $\mathbf{SL}_E$.}

Distinguons deux cas:

 a) Le groupe $S$ est de type $\mathbf A_n$ (au sens de la classification des groupes algébriques simples), donc quotient de $\mathbf{SL}(n+1)$ par un sous-groupe de son centre. Par hypothèse, $h$ divise $n+1$. D'où $\dim.S\geq\dim.\mathbf{SL}_E$, et, comme le noyau de $\varrho$ est fini, on en conclut d'abord que $\varrho$ applique $S$ sur $\mathbf{SL}_E$, puis que c'est un isomorphisme ($\mathbf{SL}_E$ étant simplement connexe).

 b) Le groupe $S$ n'est pas de type $\mathbf A_n$. La classification des groupes simples montre alors que le centre de $S$ a au plus 4 éléments. D'où $h\leq4$. Le cas $h=1$ est impossible. Le cas $h=2$ donnerait $S=\mathbf{SL}(2)$, qui est exclu.
\margin{11}
